\chapter{Traversal Law and Operational Velocity}
\label{chap:operational-velocity}

The two-speed comparison in Part V assumed constant test speeds to isolate the
\(v^2\) dependency. Operation is more general. A traversal of the running
Alignment is a monotone function \(s=s(t)\) over the interval where the chosen
orientation and applicability conditions hold.

\begin{principle}[Velocity-runtime principle]
\label{princ:velocity-runtime}
Velocity is a qualified derivative of a traversal law, not an intrinsic
Alignment property.
\end{principle}

\section{Space Becomes Time}

\begin{definition}[Velocity relation]
\label{def:velocity-relation}
For a differentiable traversal,
\[
v(t)=\dot s(t)=\frac{\dd s}{\dd t},
\qquad
a_s(t)=\dot v(t)=\frac{\dd^2s}{\dd t^2}.
\]
For a differentiable spatial quantity \(f(s)\),
\[
\frac{\dd}{\dd t}f(s(t))=v(t)f'(s(t)).
\]
\end{definition}

Applied to the separate curvature and cant laws,
\[
\frac{\dd\kappa}{\dd t}=v\kappa'(s),
\qquad
\frac{\dd u}{\dd t}=v u'(s).
\]
Their different partitions and provenance remain attached. Common traversal
does not turn them into one law or authorize a coupling rule.

For curvature acceleration, the full chain rule is
\begin{equation}
\frac{\dd^2\kappa}{\dd t^2}
=v^2\kappa''(s)+\dot v\,\kappa'(s).
\label{eq:curvature-second-time-derivative}
\end{equation}
The second term vanishes only under a stated constant-speed assumption.
Spatial smoothness belongs to \(\kappa(s)\); temporal excitation also depends
on \(s(t)\).

\section{Which Velocity Is Meant?}

\begin{definition}[Design velocity]
\label{def:design-velocity}
A design velocity is an assumption selected for a stated design purpose,
source, route or element scope, and validity.
\end{definition}

\begin{definition}[Operational velocity]
\label{def:operational-velocity}
An operational velocity is a planned, permitted, or actual traversal input
whose role and qualifications are stated for the calculation.
\end{definition}

The following roles must not be collapsed:
\begin{description}
\item[Design speed] an assumption used to develop or check a design.
\item[Permitted speed] a rule- or authority-backed upper bound under its
applicability and validity.
\item[Planned operating speed] a timetable or operating-plan input.
\item[Actual measured speed] a time-indexed observation with acquisition
provenance and uncertainty.
\item[Optimization variable] a declared degree of freedom within an
experiment, not an authorized operating value.
\item[Constant test speed] an explicit simplification used to isolate a
dependency.
\end{description}

\begin{definition}[Velocity envelope]
\label{def:velocity-envelope}
A velocity envelope is a qualified set of permitted or investigated traversal
values along an identified route or Alignment. Its source, scope, validity,
and role determine how it may be used.
\end{definition}

\section{Compact Variable-Speed Example}

At the same transition position, suppose
\[
\kappa'=2.0\times10^{-5}\,\mathrm{m^{-2}},
\qquad
\kappa''=-1.0\times10^{-7}\,\mathrm{m^{-3}}.
\]
At \(v=20\,\mathrm{m/s}\),
\[
\dot\kappa=4.0\times10^{-4}\,\mathrm{m^{-1}s^{-1}}.
\]
If speed is constant, then
\(\ddot\kappa=v^2\kappa''=-4.0\times10^{-5}\,
\mathrm{m^{-1}s^{-2}}\). If instead
\(\dot v=0.5\,\mathrm{m/s^2}\), \cref{eq:curvature-second-time-derivative}
gives
\[
\ddot\kappa
=-4.0\times10^{-5}+1.0\times10^{-5}
=-3.0\times10^{-5}\,\mathrm{m^{-1}s^{-2}}.
\]
The Alignment law is identical in both cases. Acceleration changes temporal
exposure by adding a term that a constant-speed simplification omits.

\begin{summary}
Every use of velocity must disclose source, applicability, validity, and
calculation role. The symbol \(v\) is not sufficient qualification by itself.
\end{summary}
